\documentclass{amsart}
\usepackage{amsmath,amssymb,amsthm,hyperref}
\usepackage{algorithm}
\usepackage{algpseudocode}

\newtheorem{theorem}{Theorem}
\newtheorem{lemma}{Lemma}
\newtheorem{proposition}{Proposition}
\newtheorem{corollary}{Corollary}
\theoremstyle{definition}
\newtheorem{definition}{Definition}
\newtheorem{remark}{Remark}

\newcommand{\asymp}{\sim}

\begin{document}

\title{Asymptotics of the interface width for an expanding Family--Vicsek interface}
\author{}
\date{}
\maketitle

\section{Formalization and standing assumptions}

We make precise the (semi-empirical) scaling hypotheses of the problem, state them
as axioms, and derive everything rigorously from those axioms. The mathematical
content is the analysis of a one-parameter family of competing power laws; the
experimental content is the protocol of Section~\ref{sec:protocol}.

Throughout, ``$g(t)\asymp h(t)$ as $t\to\infty$'' means there are constants
$0<c\le C<\infty$ and $t_0$ with $c\,h(t)\le g(t)\le C\,h(t)$ for all $t\ge t_0$;
for two power laws this is equivalent to equality of exponents (with the prefactor
absorbed into $c,C$), which is what we track. When we wish to keep track of the
\emph{dependence of the prefactor on a parameter} (such as the initial size $L_0$),
we say so explicitly and write, e.g., ``$g\asymp \kappa(L_0)\,t^{\theta}$ with
$\kappa(L_0)\asymp L_0^{2\alpha}$.'' We write $f(u)\asymp k(u)$ as
$u\to0^+$ or $u\to\infty$ analogously.

\begin{definition}[Family--Vicsek scaling on a fixed substrate]\label{def:FV}
Fix exponents $\alpha>0$ (roughness) and $\beta>0$ (growth), and set
$z:=\alpha/\beta>0$ (dynamic exponent). On a substrate of fixed linear size
$L_0$ the height field $h(x,t)$ has interface variance
\begin{equation}\label{eq:FV}
  W_0^2(L_0,t):=\langle h(x,t)^2\rangle-\langle h(x,t)\rangle^2
  \;\asymp\; L_0^{2\alpha}\, f\!\left(\frac{t}{L_0^{\,z}}\right),
\end{equation}
where the dimensionless positive scaling function $f$ obeys
\begin{equation}\label{eq:fasymp}
  f(u)\asymp u^{2\beta}\ \ (u\to 0^+),\qquad
  f(u)\to f_\infty\in(0,\infty)\ \ (u\to\infty),
\end{equation}
and is continuous and strictly positive on $(0,\infty)$, so that
$0<\inf_{[a,b]}f\le\sup_{[a,b]}f<\infty$ on every compact $[a,b]\subset(0,\infty)$.
\end{definition}

The single dimensionless control variable in \eqref{eq:FV} is the ratio of lateral
correlation length to system size. Defining the \emph{correlation length}
\begin{equation}\label{eq:xi}
  \xi(t):=t^{1/z},
\end{equation}
the argument of $f$ is $t/L_0^{z}=(\xi(t)/L_0)^{z}$, with crossover at
$\xi(t)=L_0$.

\subsection{The expanding interface: the problem's two regime laws as axioms}

The substrate grows as $L(t)\asymp t^{\gamma}$ ($\gamma>0$) for $t\ge t_0$, with
reference (``birth'') size $L_0:=L(t_0)$ at a fixed reference time $t_0>0$; thus
\begin{equation}\label{eq:Lt}
  L(t)=L_0\,(t/t_0)^{\gamma},\qquad t\ge t_0,\ \gamma>0.
\end{equation}
The problem \emph{posits} (``assume that \dots\ $W^2(t)$ obeys'') the following
long-time laws for the expanding-interface variance
$W^2(t):=\langle h^2\rangle-\langle h\rangle^2$:
\begin{align}
  \gamma>1/z:\quad & W^2(t)\;\asymp\; L_0^{2\alpha}\,t^{2\beta}
     \;\asymp\; L(t)^{2\alpha}\,t^{\,2\beta-2\gamma\alpha},\label{eq:axfast}\\
  \gamma<1/z:\quad & W^2(t)\;\asymp\; L_0^{2\alpha}\,t^{2\gamma\alpha}
     \;\asymp\; L(t)^{2\alpha}.\label{eq:axslow}
\end{align}
We take \eqref{eq:axfast}--\eqref{eq:axslow} as the \emph{axioms} for the expanding
interface; they are the hypotheses the problem asks us to analyze (``prove or
disprove that the two regimes above are the only asymptotic behaviors''), and the
extraction protocol of part (a) is to be derived from them. The next lemma records
that these axioms are internally consistent.

\begin{lemma}[Internal consistency of the axioms]\label{lem:consistent}
With $L(t)=L_0(t/t_0)^{\gamma}$ the two displayed forms in \eqref{eq:axfast} are
equal up to the constant factor $t_0^{2\gamma\alpha}$, and likewise the two forms
in \eqref{eq:axslow} are equal up to $t_0^{2\gamma\alpha}$. In particular each
regime law is a single well-defined power law in $t$ whose prefactor scales as
$L_0^{2\alpha}$:
\[
   W^2(t)\asymp \kappa(L_0)\,t^{\theta(\gamma)},\qquad
   \kappa(L_0)\asymp L_0^{2\alpha},\qquad
   \theta(\gamma)=\begin{cases}2\beta,&\gamma>1/z,\\ 2\gamma\alpha,&\gamma<1/z.\end{cases}
\]
\end{lemma}
\begin{proof}
Using \eqref{eq:Lt}, $L(t)^{2\alpha}=L_0^{2\alpha}t_0^{-2\gamma\alpha}t^{2\gamma\alpha}$.
Hence in \eqref{eq:axfast},
$L(t)^{2\alpha}t^{2\beta-2\gamma\alpha}
=L_0^{2\alpha}t_0^{-2\gamma\alpha}\,t^{2\beta}$,
which equals $L_0^{2\alpha}t^{2\beta}$ up to the constant $t_0^{-2\gamma\alpha}$.
In \eqref{eq:axslow}, $L(t)^{2\alpha}=L_0^{2\alpha}t_0^{-2\gamma\alpha}t^{2\gamma\alpha}$
equals $L_0^{2\alpha}t^{2\gamma\alpha}$ up to the same constant. In both regimes the
$t$-power is unique and the $L_0$-prefactor is $\asymp L_0^{2\alpha}$.
\end{proof}

\begin{remark}[Relation to the standard adiabatic Family--Vicsek closure]
\label{rem:closure}
The axioms \eqref{eq:axfast}--\eqref{eq:axslow} are \emph{very close to}, but not
identical with, what one gets from the naive ``adiabatic'' closure
\begin{equation}\label{eq:adiabatic}
  W^2(t)\;\overset{?}{\asymp}\;L(t)^{2\alpha}f\!\Big(\tfrac{t}{L(t)^{z}}\Big),
\end{equation}
in which the \emph{fixed} size in \eqref{eq:FV} is replaced at each instant by the
current size $L(t)$. A direct computation (Lemma~\ref{lem:adiab} below) shows that
\eqref{eq:adiabatic} reproduces the \emph{slow} law \eqref{eq:axslow} exactly, but
in the \emph{fast} regime it gives $W^2\asymp t^{2\beta}$ with an
\emph{$L_0$-independent} prefactor, i.e.\ it \emph{omits} the factor $L_0^{2\alpha}$
present in \eqref{eq:axfast}. The two models differ by exactly one physical
assumption about the growth (unsaturated) phase:
\begin{itemize}
\item \textbf{Adiabatic closure \eqref{eq:adiabatic}:} the lateral amplitude is
  continually \emph{reset} to the current size; the vertical width during growth is
  $t^{2\beta}$, with no memory of the birth size $L_0$.
\item \textbf{Problem's axiom \eqref{eq:axfast}:} the lateral amplitude is
  \emph{set by the birth substrate} $L_0$ (the substrate on which the interface was
  seeded), so the same temporal growth $t^{2\beta}$ is multiplied by the
  birth-amplitude $L_0^{2\alpha}$.
\end{itemize}
Both are legitimate phenomenological closures; the problem selects the second. We
therefore take \eqref{eq:axfast}--\eqref{eq:axslow} at face value as the governing
hypotheses (this is essential for part (a), which asks one to read \emph{both}
$\alpha$ and $\beta$ off the fast regime via the $L_0$- and $t$-dependence), and we
record the adiabatic computation only to delineate precisely where the two models
agree and differ. \emph{All results below are theorems conditional on the stated
axioms.}
\end{remark}

\begin{lemma}[The adiabatic closure: where it agrees and differs]\label{lem:adiab}
Assume \eqref{eq:FV}--\eqref{eq:fasymp} and \eqref{eq:adiabatic}, with
$L(t)=L_0(t/t_0)^{\gamma}$. Put $\rho(t):=t/L(t)^{z}$. Then
$\rho(t)\asymp t^{z(1/z-\gamma)}$, and:
\begin{itemize}
\item if $\gamma<1/z$ then $\rho(t)\to\infty$ and
  $W^2\asymp f_\infty L(t)^{2\alpha}\asymp L_0^{2\alpha}t^{2\gamma\alpha}$,
  agreeing with \eqref{eq:axslow};
\item if $\gamma>1/z$ then $\rho(t)\to0$ and
  $W^2\asymp L(t)^{2\alpha}\rho(t)^{2\beta}=t^{2\beta}$, with $L_0$-independent
  prefactor (since $2\alpha-2\beta z=0$), differing from \eqref{eq:axfast} by the
  factor $L_0^{2\alpha}$.
\end{itemize}
\end{lemma}
\begin{proof}
By \eqref{eq:Lt},
$\rho(t)=t/(L_0(t/t_0)^{\gamma})^{z}=L_0^{-z}t_0^{\gamma z}t^{1-\gamma z}
\asymp t^{z(1/z-\gamma)}$ (as $z>0$). If $\gamma<1/z$ the exponent $1-\gamma z>0$ so
$\rho\to\infty$; by \eqref{eq:fasymp} $f(\rho)\to f_\infty$, and
$W^2\asymp f_\infty L(t)^{2\alpha}=f_\infty L_0^{2\alpha}t_0^{-2\gamma\alpha}t^{2\gamma\alpha}$.
If $\gamma>1/z$ the exponent is negative, $\rho\to0$; by \eqref{eq:fasymp}
$f(\rho)\asymp\rho^{2\beta}$ for large $t$, and
$W^2\asymp L(t)^{2\alpha}\rho^{2\beta}=L(t)^{2\alpha}(t/L(t)^z)^{2\beta}
=t^{2\beta}L(t)^{2\alpha-2\beta z}=t^{2\beta}$ since $2\alpha-2\beta z
=2\alpha-2\beta(\alpha/\beta)=0$, with prefactor independent of $L_0$.
\end{proof}

\section{The dichotomy theorem: exactly two leading regimes (plus the boundary)}
\label{sec:dichotomy}

We now work \emph{from the axioms} \eqref{eq:axfast}--\eqref{eq:axslow}, extended to
the boundary $\gamma=1/z$ by continuity of the underlying Family--Vicsek crossover.
To do this rigorously at the boundary we make the (weakest possible) hypothesis that
the expanding-interface law is the adiabatic crossover whenever the argument of $f$
stays in a compact set; this fixes the boundary case and does not affect the two
strict regimes, which are governed by \eqref{eq:axfast}--\eqref{eq:axslow}.

\begin{theorem}[Trichotomy of asymptotics]\label{thm:main}
Assume \eqref{eq:axfast}--\eqref{eq:axslow}, with the boundary case governed by the
adiabatic crossover \eqref{eq:adiabatic}. As $t\to\infty$, with
$L(t)=L_0(t/t_0)^{\gamma}$,
\begin{align}
  \text{\emph{(F)} fast/uncorrelated, } \gamma>1/z:\quad
    & W^2(t)\;\asymp\; L_0^{2\alpha}\,t^{2\beta}
      \;\asymp\; L(t)^{2\alpha}\,t^{\,2\beta-2\gamma\alpha},
      \label{eq:fast}\\[2pt]
  \text{\emph{(S)} slow/correlated, } \gamma<1/z:\quad
    & W^2(t)\;\asymp\; L_0^{2\alpha}\,t^{2\gamma\alpha}
      \;\asymp\; L(t)^{2\alpha},
      \label{eq:slow}\\[2pt]
  \text{\emph{(C)} critical, } \gamma=1/z:\quad
    & W^2(t)\;\asymp\; t^{2\beta}\;\asymp\;L(t)^{2\alpha}.
      \label{eq:crit}
\end{align}
In particular, in each strict regime $W^2(t)\asymp \kappa(L_0)\,t^{\theta(\gamma)}$
for a single power
\begin{equation}\label{eq:theta}
  \theta(\gamma)=\min\{\,2\beta,\;2\gamma\alpha\,\}
  =\begin{cases}2\beta, & \gamma\ge 1/z,\\[2pt] 2\gamma\alpha, & \gamma\le 1/z,\end{cases}
\end{equation}
and there are \emph{no other} asymptotic power laws. The amplitude scales as
\begin{equation}\label{eq:kappa}
  \kappa(L_0)\asymp
  \begin{cases}
    L_0^{2\alpha}, & \gamma>1/z \quad(\text{regime (F)}),\\[2pt]
    L_0^{2\alpha}, & \gamma<1/z \quad(\text{regime (S)}).
  \end{cases}
\end{equation}
The exponent $\theta(\gamma)$ is continuous, non-decreasing, concave (piecewise
linear with a single kink at $\gamma=1/z$), and equals $2\beta$ on $[1/z,\infty)$.
\end{theorem}

\begin{proof}
\textbf{Case (F): $\gamma>1/z$.} This is axiom \eqref{eq:axfast}. By
Lemma~\ref{lem:consistent}, the two displayed forms agree and the $t$-power is
$2\beta$ with prefactor $\asymp L_0^{2\alpha}$. Since
$\gamma>1/z=\beta/\alpha$ gives $2\gamma\alpha>2\beta$, the exponent is
$2\beta=\min\{2\beta,2\gamma\alpha\}$, matching \eqref{eq:theta}.

\medskip\noindent\textbf{Case (S): $\gamma<1/z$.} This is axiom \eqref{eq:axslow}.
By Lemma~\ref{lem:consistent}, the $t$-power is $2\gamma\alpha$ with prefactor
$\asymp L_0^{2\alpha}$. Since $\gamma<\beta/\alpha$ gives $2\gamma\alpha<2\beta$,
the exponent is $2\gamma\alpha=\min\{2\beta,2\gamma\alpha\}$.

\medskip\noindent\textbf{Case (C): $\gamma=1/z$.} By Lemma~\ref{lem:adiab}'s
computation, $\rho(t)=t/L(t)^z\asymp t^{z(1/z-\gamma)}=t^0\asymp1$, so
$\rho(t)\in[a,b]$ for some $0<a\le b<\infty$ and all large $t$. By continuity and
positivity of $f$ on the compact $[a,b]$ (Definition~\ref{def:FV}),
$f(\rho(t))\asymp1$, hence by \eqref{eq:adiabatic}
$W^2(t)\asymp L(t)^{2\alpha}\asymp t^{2\gamma\alpha}=t^{2\alpha/z}=t^{2\beta}$,
using $\gamma=1/z$ and $z=\alpha/\beta$. This is \eqref{eq:crit}.

\medskip\noindent\textbf{No other behaviour; shape of $\theta$.} The cases
partition $(0,\infty)$; in each, $W^2$ is a single power $t^{\theta(\gamma)}$ with
\eqref{eq:theta}. Continuity at $\gamma=1/z$: both branches give
$2\beta=2(1/z)\alpha$. On $(0,1/z]$, $\theta=2\alpha\gamma$ strictly increases
(slope $2\alpha$); on $[1/z,\infty)$, $\theta\equiv2\beta$. Being the minimum of
two affine functions of $\gamma$, $\theta$ is concave with a single downward kink at
$\gamma=1/z$ (slope drops $2\alpha\to0$). This disproves any third leading regime:
the only freedom in $\theta$ is the kink location $1/z$ and the slopes $2\alpha,0$.
\end{proof}

\begin{remark}[On the literal reading of the problem's fast form]\label{rem:correction}
The problem's two fast-regime expressions,
$L_0^{2\alpha}t^{2\beta}$ and $L(t)^{2\alpha}t^{2\beta-2\gamma\alpha}$, are
\emph{mutually consistent}: by Lemma~\ref{lem:consistent} they differ only by the
constant $t_0^{2\gamma\alpha}$, and \emph{both} carry the birth-size prefactor
$L_0^{2\alpha}$. (It would be an error to claim the second form equals $t^{2\beta}$;
it equals $L_0^{2\alpha}t_0^{-2\gamma\alpha}t^{2\beta}$.) The only subtlety is the
relation to the naive adiabatic closure \eqref{eq:adiabatic}, which would
\emph{drop} the $L_0^{2\alpha}$ factor in the fast regime (Lemma~\ref{lem:adiab});
this is a genuine modeling choice, not an inconsistency, and the problem fixes it by
fiat (the lateral amplitude is set by the birth substrate). Consequently $\alpha$
\emph{can} be read from the $L_0$-dependence of the fast regime, exactly as part (a)
requests; see Section~\ref{sec:protocol}. \emph{Caveat for experimentalists:} if the
true system follows the adiabatic closure rather than \eqref{eq:axfast}, the fast
regime would be $L_0$-independent and $\alpha$ would instead have to be taken from
the slow regime; the protocol below includes a cross-check (Step~3 / regime route)
that detects which closure holds.
\end{remark}

\begin{remark}[The boundary case]
The two strict regimes leave $\gamma=1/z$ to the crossover; there $\rho(t)\asymp1$,
so $f$ sits in its crossover, a genuinely distinct sub-leading structure. But the
leading power $t^{2\beta}$ at $\gamma=1/z$ matches both one-sided limits, so
$\theta(\gamma)$ is continuous and the two strict regimes are the only two
\emph{leading} power laws, separated by the single critical point $\gamma=1/z$.
\end{remark}

\section{The order parameter $1/z-\gamma$: distance from correlation}
\label{sec:correlation}

\begin{definition}[Global correlation]\label{def:corr}
At time $t$ the interface is \emph{globally correlated} (saturated) if
$\xi(t)\ge L(t)$ and \emph{uncorrelated} (roughening) if $\xi(t)<L(t)$.
\end{definition}

\begin{proposition}[Sign law]\label{prop:sign}
Let $\Delta:=1/z-\gamma$. For large $t$,
\[
  \frac{\xi(t)}{L(t)}\asymp t^{\Delta},\qquad
  \frac{\log(\xi(t)/L(t))}{\log t}\to\Delta .
\]
Hence: (i) $\Delta>0\iff\gamma<1/z\iff\xi/L\to\infty$: correlated phase (S);
(ii) $\Delta<0\iff\gamma>1/z\iff\xi/L\to0$: uncorrelated phase (F);
(iii) $\Delta=0\iff\gamma=1/z\iff\xi/L\asymp1$: critical phase (C).
The magnitude $|\Delta|$ is the per-decade rate of approach to/retreat from full
correlation: $\xi/L$ changes by the factor $10^{\Delta}$ when $t\mapsto10t$. Thus
$\Delta=1/z-\gamma$ is the precise ``distance from correlation,'' its sign the phase
and $|\Delta|$ the speed.
\end{proposition}
\begin{proof}
By \eqref{eq:xi},\eqref{eq:Lt}, $\xi/L=t^{1/z}/(L_0(t/t_0)^{\gamma})
\asymp t^{1/z-\gamma}=t^{\Delta}$, giving the limit. The equivalences are the sign
cases of $\Delta$, matched to the cases of Theorem~\ref{thm:main} via
Lemma~\ref{lem:adiab} ($\rho=(\xi/L)^z\asymp t^{z\Delta}$, same sign). Per decade:
$(\xi/L)(10t)\big/(\xi/L)(t)\asymp10^{\Delta}$.
\end{proof}

\begin{theorem}[Minimal control strength to switch phase]\label{thm:control}
Let the bare system have growth index $\gamma_0$ and dynamic exponent $z$, hence
phase $\operatorname{sgn}(\Delta_0)$ with $\Delta_0:=1/z-\gamma_0$. A control
protocol replacing $\gamma_0$ by $\gamma=\gamma_0+\delta$ ($\delta\in\mathbb R$,
$\gamma>0$) crosses the phase boundary iff $\gamma$ crosses $1/z$, i.e.\ iff
$\delta$ crosses $\Delta_0$. The boundary is reached exactly at $\delta=\Delta_0$
(critical phase (C)), and the system lands strictly in the opposite phase for any
$\delta$ on the far side of $\Delta_0$; the infimal switching strength is
\begin{equation}\label{eq:thresh}
  |\delta_{\min}|=\big|\,\tfrac1z-\gamma_0\,\big|=|\Delta_0|,
\end{equation}
attained in the limit. Directionally: from the uncorrelated phase
($\gamma_0>1/z$) one must \emph{decrease} $\gamma$ with
$\delta<\Delta_0<0$, $|\delta|>|\Delta_0|$; from the correlated phase
($\gamma_0<1/z$) one must \emph{increase} $\gamma$ with $\delta>\Delta_0>0$.
\end{theorem}
\begin{proof}
By Theorem~\ref{thm:main} and Proposition~\ref{prop:sign}, the phase is
$\operatorname{sgn}(1/z-\gamma)=\operatorname{sgn}(\Delta)$, boundary $\Delta=0$.
With $\gamma=\gamma_0+\delta$, $\Delta(\delta)=\Delta_0-\delta$ is affine, strictly
decreasing, with unique zero $\delta=\Delta_0$. Thus the sign flips exactly as
$\delta$ passes $\Delta_0$: same side $\Rightarrow$ same phase, far side
$\Rightarrow$ opposite phase, $\delta=\Delta_0\Rightarrow$ critical. The infimal
$|\delta|$ that changes the sign is $|\Delta_0|$ (not attained at the strict level,
attained as the boundary). The directional statements are the two sign cases of
$\Delta_0$.
\end{proof}

\section{Experimental protocol and unambiguous extraction of $\alpha,\beta$}
\label{sec:protocol}

By Theorem~\ref{thm:main} (axioms \eqref{eq:axfast}--\eqref{eq:axslow}), in the
\textbf{fast} regime (F) the variance carries \emph{both} exponents:
\begin{equation}\label{eq:fastlog}
  \log W^2(t)\;=\;2\alpha\,\log L_0\;+\;2\beta\,\log t\;+\;\text{const}\;+\;o(\log\,\cdot),
\end{equation}
so $\beta$ is the time-slope and $\alpha$ is half the $L_0$-slope. This is exactly
the situation part (a) describes (``estimate $\alpha,\beta$ \dots\ from the
long-time dependence of $W^2(t)$ on the initial substrate size $L_0$ and on time
$t$ \dots\ in the fast-growth regime $\gamma>1/z$''). The \textbf{slow} regime (S)
provides an independent cross-check, since there
\begin{equation}\label{eq:slowlog}
  \log W^2(t)\;=\;2\alpha\,\log L_0\;+\;2\gamma\alpha\,\log t\;+\;\text{const}\;+\;o(\log\,\cdot),
\end{equation}
which determines $\alpha$ from the $L_0$-slope and $\gamma\alpha$ from the
time-slope.

\subsection*{Protocol (biological realization)}
\begin{enumerate}
\item Prepare $m\ge2$ replicate colonies/tissues with controlled distinct initial
  front sizes $L_0^{(1)}<\cdots<L_0^{(m)}$, all other conditions identical so that
  $\alpha,\beta,z,f$ are shared.
\item Operate the system in the fast-growth setting $\gamma>1/z$ (the natural front
  expansion of a spreading colony typically realizes $\gamma>1/z$; if not, increase
  the expansion rate by an external control, see below). Confirm $\gamma>1/z$
  a~posteriori from the fitted exponents (Step~6 of the algorithm).
\item For each replicate and a set of large times $t_1<\cdots<t_n$ (in the
  asymptotic window $t\gg t_0$), image the front, extract $h(x,t)$, and compute the
  spatial variance $W^2(t)=\overline{h^2}-\overline h^2$ (replicate-averaged).
\item Track the front size $L(t)$ and fit $\log L\sim\gamma\log t$ to measure
  $\gamma$.
\item (Optional cross-check.) Run an additional slow setting $\gamma<1/z$ to verify
  $\alpha$ via \eqref{eq:slowlog}; Theorem~\ref{thm:control} gives the minimal
  change $|\delta|>|1/z-\gamma|$ needed to move between the regimes.
\end{enumerate}

\subsection*{Extraction algorithm}

\begin{algorithm}[H]
\caption{Unambiguous extraction of $\alpha,\beta,z$ from the fast regime}
\begin{algorithmic}[1]
\Require Variances $W^2(L_0^{(i)},t_j)$ in a fast setting $\gamma>1/z$ for $m\ge2$
  initial sizes and $n\ge2$ large times; measured $\gamma$.
\Ensure $\hat\beta,\hat\alpha,\hat z$ with self-consistency checks.
\State \textbf{($\beta$ from the time direction.)} For each fixed $L_0^{(i)}$,
  regress $\log W^2(L_0^{(i)},t_j)$ on $\log t_j$; the slope is $2\hat\beta_i$.
  Check $\hat\beta_1\approx\cdots\approx\hat\beta_m$ (the time exponent must be
  $L_0$-independent) and set $\hat\beta\gets\tfrac12\,\mathrm{mean}_i(2\hat\beta_i)$.
\State \textbf{($\alpha$ from the size direction.)} For each fixed $t_j$, regress
  $\log W^2(L_0^{(i)},t_j)$ on $\log L_0^{(i)}$ across $i=1,\dots,m$; the slope is
  $2\hat\alpha_j$. Check $\hat\alpha_1\approx\cdots\approx\hat\alpha_n$
  (time-independence of the size exponent) and set
  $\hat\alpha\gets\tfrac12\,\mathrm{mean}_j(2\hat\alpha_j)$.
\State \textbf{(Joint fit / cross-check.)} Fit the bilinear model
  $\log W^2=a+2\alpha\log L_0+2\beta\log t$ by least squares over all $mn$ points;
  the fitted $(\alpha,\beta)$ must agree with Steps~1--2. (If a slow-setting dataset
  is available, also fit \eqref{eq:slowlog} and check the $\alpha$ values agree;
  disagreement of the fast $L_0$-slope with the slow $L_0$-slope flags an adiabatic
  closure, cf.\ Remark~\ref{rem:correction}.)
\State \textbf{(Dynamic exponent.)} $\hat z\gets\hat\alpha/\hat\beta$.
\State \textbf{(Regime check.)} Verify $\gamma>1/\hat z$ (consistency of the
  fast-regime assumption).
\State \Return $\hat\alpha,\hat\beta,\hat z$.
\end{algorithmic}
\end{algorithm}

\begin{proposition}[Correctness, identifiability, termination, complexity]\label{prop:alg}
Under \eqref{eq:axfast}--\eqref{eq:axslow}:
\begin{enumerate}
\item[\textnormal{(i)}] \emph{(Consistency.)} In the fast setting
  ($\gamma>1/z$), \eqref{eq:fastlog} holds, so the Step-1 time slope $\to2\beta$ and
  the Step-2 size slope $\to2\alpha$; hence $\hat\beta\to\beta$,
  $\hat\alpha\to\alpha$, $\hat z\to z$.
\item[\textnormal{(ii)}] \emph{(Unambiguous identifiability.)} In the fast regime,
  $(\log L_0,\log t)\mapsto\log W^2$ is affine with gradient $(2\alpha,2\beta)$. The
  two regressors $\log L_0$ (varied over $m\ge2$ distinct sizes) and $\log t$
  (varied over $n\ge2$ times) are independently controlled, so the $(\le mn)\times3$
  design matrix $\Phi$ with rows $(1,\log L_0^{(i)},\log t_j)$ has full column rank
  $3$; therefore the coefficient vector $(\text{const},2\alpha,2\beta)$ is uniquely
  determined and $(\alpha,\beta)$ — hence $z=\alpha/\beta$ — are uniquely fixed.
\item[\textnormal{(iii)}] \emph{(The fast regime suffices.)} Because
  \eqref{eq:axfast} carries both $L_0^{2\alpha}$ and $t^{2\beta}$, the fast regime
  alone identifies both exponents, answering part (a) as posed; the slow regime
  \eqref{eq:axslow} provides an independent determination of $\alpha$ (and of
  $\gamma\alpha$) for cross-validation, and serves to discriminate the present model
  from the adiabatic closure (Remark~\ref{rem:correction}).
\item[\textnormal{(iv)}] \emph{(Termination/complexity.)} The algorithm performs
  $O(1)$ regressions over $\le mn$ points; it runs in $O(mn)$ time and $O(mn)$
  space and halts (loops over finite index sets only).
\end{enumerate}
\end{proposition}
\begin{proof}
(i) Taking logs of \eqref{eq:fast} gives
$\log W^2=2\alpha\log L_0+2\beta\log t+\text{const}+r(L_0,t)$, where the residual
$r=\log\bigl(W^2/(\text{leading power})\bigr)$ is, by the definition of $\asymp$,
\emph{bounded}: $\log c\le r\le\log C$ for all large $t$ (with $c,C$ the
$\asymp$-constants). We use the elementary OLS bound: for data $y_k=s\,x_k+a+r_k$
with $|r_k|\le R$ over points $x_k$ ($k=1,\dots,N$), the OLS slope is
$\hat s=s+\dfrac{\sum_k(x_k-\bar x)r_k}{\sum_k(x_k-\bar x)^2}$, and by
Cauchy--Schwarz
\[
   |\hat s-s|\le\frac{R\sqrt N\,\sqrt{\sum_k(x_k-\bar x)^2}}{\sum_k(x_k-\bar x)^2}
   =\frac{R}{\sqrt{\operatorname{Var}_k(x_k)}}\xrightarrow[\operatorname{Var}\to\infty]{}0 .
\]
Thus the Step-1 time slope (with $x_k=\log t_j$ spread over many decades) converges
to $2\beta$ and the Step-2 size slope (with $x_k=\log L_0^{(i)}$ spread over several
decades) converges to $2\alpha$. Hence $\hat\beta\to\beta$, $\hat\alpha\to\alpha$,
$\hat z\to z$. (When the residual is exactly constant — pure power law — $m=2$ sizes
and $n=2$ times already recover the slopes exactly.)
(ii) Stack the fast data as $y=\Phi w$, $w=(\text{const},2\alpha,2\beta)^\top$, rows
$(1,\log L_0^{(i)},\log t_j)$. With $m\ge2$ distinct $L_0$ and $n\ge2$ distinct $t$,
the three columns (constant; nonconstant function of $i$ alone; nonconstant function
of $j$ alone) are linearly independent, so $\Phi$ has rank $3$ and $w$ — hence
$(\alpha,\beta)$ — is unique.
(iii) Is the content of \eqref{eq:fast}: both exponents appear in the fast law, and
\eqref{eq:slow} supplies $\alpha$ and $\gamma\alpha$ independently.
(iv) Immediate from the operation count over finite index sets.
\end{proof}

\begin{remark}[Closure of part (b)]
Proposition~\ref{prop:alg} yields $\hat z=\hat\alpha/\hat\beta$, and with the
measured $\gamma$ the order parameter $\widehat\Delta=1/\hat z-\gamma$ of
Proposition~\ref{prop:sign} classifies the phase; by Theorem~\ref{thm:control} the
minimal control strength to switch phase is
$|\delta_{\min}|=|1/\hat z-\gamma|=|\widehat\Delta|$.
\end{remark}

\section{Summary}

\begin{itemize}
\item \textbf{Trichotomy.} Under the problem's regime axioms
  \eqref{eq:axfast}--\eqref{eq:axslow}, the long-time width is the single power law
  $W^2(t)\asymp\kappa(L_0)\,t^{\theta(\gamma)}$,
  $\theta(\gamma)=\min\{2\beta,2\alpha\gamma\}$, $\kappa(L_0)\asymp L_0^{2\alpha}$
  in both regimes (Theorem~\ref{thm:main}). The two announced regimes (F)
  $\gamma>1/z$ and (S) $\gamma<1/z$ are exactly the two leading behaviours,
  separated by the single critical point $\gamma=1/z$ (regime (C)); there are no
  others. This settles ``prove or disprove'': the two strict regimes are the only
  leading power laws, and the boundary case (omitted in the problem) is identified.
\item \textbf{Consistency, not correction.} The problem's two fast-regime forms
  $L_0^{2\alpha}t^{2\beta}$ and $L(t)^{2\alpha}t^{2\beta-2\gamma\alpha}$ are
  \emph{mutually consistent} (equal up to a $t_0$-constant) and both carry
  $L_0^{2\alpha}$ (Lemma~\ref{lem:consistent}); the only nuance is that the
  naive adiabatic closure \eqref{eq:adiabatic} would omit this $L_0^{2\alpha}$ in
  the fast phase (Lemma~\ref{lem:adiab}), a distinct modeling choice the problem
  does not adopt (Remark~\ref{rem:correction}).
\item \textbf{Order parameter \& control.} $\Delta=1/z-\gamma$ equals
  $\lim_t\log(\xi/L)/\log t$; its sign is the phase and $|\Delta|$ the per-decade
  rate (Proposition~\ref{prop:sign}). Minimal control to switch phase is
  $|\delta_{\min}|=|1/z-\gamma_0|$ (Theorem~\ref{thm:control}).
\item \textbf{Extraction (part (a)).} In the fast regime both exponents are
  identifiable: $\beta$ from the time-slope and $\alpha$ from the $L_0$-slope of
  $W^2$ (Algorithm~1, Proposition~\ref{prop:alg}); the slow regime cross-checks
  $\alpha$. The design is full rank, so $(\alpha,\beta,z)$ are unambiguously
  determined.
\end{itemize}

\end{document}
